\section{Discussion}
\label{sec:discussion}

%-------------------------------------------------------------------------
\subsection{Why the Original Diffusion Loss Fails}
\label{sec:diff_loss}
% [3 pts] Intuitively explain why L_Diff = ||eps_hat(x_t, y, t) - eps_t||^2 did not produce
% realistic samples (use < 1 math equation) & provide visual results of applying L_Diff.

For the three prompts in \cref{fig:diff_loss}, we directly minimized the noise-prediction MSE by
backpropagating through the frozen U-Net, retaining its input Jacobian.

Noise-prediction error measures how well the model recovers a known perturbation; it does not
directly measure whether the clean latent represents a realistic image.
When the latent is optimized, the U-Net Jacobian transforms the noise residual into an update.
This transformation can amplify sensitivity to local patterns rather than produce a coherent change
toward the requested object.
SDS omits that term and uses the noise residual directly as guidance.

In \cref{fig:diff_loss}, the diffusion-loss results contain dense, colorful noise without a
recognizable dog, castle, or car, whereas SDS produces recognizable, though blurred, subjects.
The average CLIP score over these three prompts is 0.1831 for diffusion loss versus 0.2931 for SDS.
The images are consistent with optimizing denoising error without obtaining coherent object
structure.

%-------------------------------------------------------------------------
\subsection{Why SDI and VSD Outperform SDS}
\label{sec:why_better}
% [3 pts] Intuitively analyze why SDI and VSD perform better than SDS (without heavy math).

\noindent\textbf{SDS.}
In \cref{fig:qualitative}, SDS produces recognizable but heavily smoothed subjects, such as the
burger and cabin, with little surface texture or fine structural detail.
SDS samples fresh Gaussian noise and a timestep at each iteration, so successive updates can favor
different image changes.
These fluctuating directions can wash out fine details during optimization.
Strong CFG reinforces the prompt, but does not ensure that these details are preserved.

\noindent\textbf{VSD.}
VSD still adds Gaussian noise, but replaces that sampled noise in the distillation residual with
the prediction of a LoRA model trained to denoise the current generated latent.
The update therefore compares the pretrained image prior with a predictor adapted to the current
samples, allowing it to account for their distribution rather than use a single sampled noise
as the reference.
In \cref{fig:qualitative}, VSD recovers sharper burger ingredients and more vegetation around the
villa than SDS.
The villa's CLIP score also improves from 0.2935 to 0.3097 (\cref{tab:clip}).

\noindent\textbf{SDI.}
SDI constructs the noisy latent by DDIM inversion of the current image instead of injecting fresh
Gaussian noise.
This makes the perturbation depend on the image and prompt, supporting more coherent updates.
Its decreasing timestep schedule moves from high to low noise, while the cached target provides
a fixed reconstruction objective between refreshes.
In \cref{fig:qualitative}, SDI produces clearer cabin and church structures and cleaner surfaces
than SDS.
Its average CLIP score increases from 0.3042 to 0.3086.